\documentclass[11pt]{article}
\usepackage{fontspec,amsmath,amssymb,amsthm,hyperref}
\usepackage[margin=1in]{geometry}
\setlength{\emergencystretch}{3em}
\title{The actual scaled particle force and its geometry}
\author{LangevinLean mathematical exposition}
\date{}
\begin{document}\maketitle
\noindent Source module: \texttt{KIPLMC/Geometry.lean}. Compile this proof document with XeLaTeX or LuaLaTeX. The accompanying HTML pairs the mathematical argument with the exact Lean statements and records the scope of the formalization.


The shared IPLA state stores \(z=(\theta,X^1/\sqrt N,\ldots,X^N/\sqrt N)\) in a Euclidean coordinate space with independent parameter and latent dimensions. This module keeps physical time, scaling velocities by the same spatial map. In this convention the kinetic temperature is \(1/N\), the potential is the particle average \(\bar U_N\), and the force is its actual gradient. It is equivalent to the manuscript's additional time change but avoids reimplementing the shared kinetic step.


\section{Define the average and prove that its gradient is the force}


\texttt{State} reuses the scaled Euclidean IPLA state. The existing lifted potential is \(V_N(z)=\sum_iU(\theta,\sqrt Nz_i)\). Define \(\bar U_N=V_N/N\) and \(F=\nabla V_N/N\). To prove \(\nabla\bar U_N=F\), compare their inner products against an arbitrary vector. The derivative rule for a scalar multiple gives \(D\bar U_N[w]=N^{-1}DV_N[w]\), and the inner-product characterization of gradients identifies the two vectors. Lean uses \texttt{ext\_inner\_right}, the proved differentiability of the lifted potential, and \texttt{HasFDerivAt.const\_mul}. For \(N>0\), cancellation of the nonzero particle number gives \(V_N=N\bar U_N\).



\noindent Lean declarations: \texttt{State}, \texttt{averagePotential}, \texttt{force}, \texttt{gradient\_averagePotential}, \texttt{potential\_eq\_mul\_average}.

\section{Average the lifted strong convexity}


The existing particle geometry proves \(V_N\) is \(Nm\)-strongly convex whenever the pair potential is \(m\)-strongly convex. Multiply its convex-combination inequality by \(N^{-1}>0\): the function terms become \(\bar U_N\), and the quadratic correction \(N^{-1}(Nm)\|x-y\|^2/2\) becomes \(m\|x-y\|^2/2\). Lean constructs the strong-convexity structure directly, obtains the lifted inequality, multiplies it by the nonnegative inverse particle number, and uses \texttt{field\_simp} and \texttt{ring} to cancel constants. The proof requires positive particle count but imposes no artificial equality of the parameter and latent dimensions.



\noindent Lean declarations: \texttt{averagePotential\_strongConvex}.

\section{Recover the mean parameter force and unaveraged latent force}


The parameter block of \(\nabla V_N\) is the sum of pair parameter gradients, so \(F_\theta=N^{-1}\sum_i\nabla_\theta U(\theta,X^i)\). Each latent block of \(\nabla V_N\) is \(\sqrt N\nabla_xU\). Dividing by \(N\) and undoing the spatial scaling gives \(\sqrt N F_i=\nabla_xU(\theta,X^i)\). The parameter proof simplifies with the existing gradient-coordinate theorem. The latent proof expands coordinate scalar multiplication, groups the coefficient into \((\sqrt N)^2N^{-1}\), rewrites \((\sqrt N)^2=N\), and cancels the nonzero \(N\). These identities verify the actual force used by both particle algorithms.



\noindent Lean declarations: \texttt{force\_parameter}, \texttt{force\_particle}.

\section{The force retains the pair smoothness constant}


The shared lifted gradient has Lipschitz constant \(NL\). Since \(F=\nabla V_N/N\), its difference satisfies \(\|F(x)-F(y)\|\le N^{-1}(NL)\|x-y\|=L\|x-y\|\). Lean starts from the existing lifted Lipschitz estimate, rewrites norms of scalar multiples with the nonnegative inverse particle count, multiplies the inequality by that inverse, and cancels \(N\). This proves that the physical-time scaled force has a smoothness constant independent of the number of particles, a key feature of the favorable KIPLMC1 scaling.



\noindent Lean declarations: \texttt{force\_lipschitz}.

\end{document}
